\documentclass[../DoAn.tex]{subfiles}
\begin{document}
\section{Kết luận}

Đồ án đã hoàn thành các mục tiêu đề ra ban đầu: thiết kế, huấn luyện và triển khai thành công một hệ thống tự động nhận dạng biển số xe Việt Nam (ALPR) dựa trên phân tích luồng video. Kết quả thực nghiệm xác nhận hệ thống có khả năng nhận dạng chính xác toàn chuỗi ký tự trên cả biển số một dòng và hai dòng, có khả năng hoạt động trên các video khảo sát.

Đóng góp thực tiễn của đồ án là việc xây dựng hoàn chỉnh một quy trình (pipeline) xử lý đa tầng chuyên biệt cho bài toán ALPR. Thay vì phụ thuộc vào quyết định trên một khung hình đơn lẻ tiềm ẩn nhiều rủi ro, hệ thống đã chuyển hướng tiếp cận sang việc theo vết phương tiện (tracking) và đối soát bằng chứng trên nhiều khung hình (multi-frame). Cách tiếp cận này giúp giảm thiểu sai sót, tăng cường độ tin cậy và phản ánh sát với cách con người quan sát dữ liệu biển số trong thực tế.

Thành quả cốt lõi của đồ án nằm ở khả năng kết hợp hiệu quả các mô hình AI thành phần thành một hệ thống đồng nhất và có khả năng chịu lỗi tốt. Đặc biệt, đồ án cho thấy hiệu quả trong phạm vi hai kịch bản thử nghiệm đối với việc chọn lọc, tinh chỉnh (fine-tuning) và huấn luyện lại các kiến trúc mạng nơ-ron nhỏ gọn. Việc sử dụng các mô hình nhỏ gọn nhưng được tối ưu hóa chuyên sâu cho đặc thù dữ liệu Việt Nam cho thấy tính khả thi của hướng đa khung hình và được thiết kế nhằm hạn chế số lần OCR không cần thiết, tránh sự lãng phí khi lạm dụng các siêu mô hình ngôn ngữ - thị giác.

Bên cạnh đóng góp về mô hình và pipeline nhận dạng, đồ án đã triển khai một sản phẩm Web end-to-end để đưa toàn bộ quy trình nhận dạng vào môi trường vận hành thử nghiệm. Hệ thống hỗ trợ tải lên video, tiếp nhận nguồn video streaming, theo dõi tiến trình xử lý, hiển thị kết quả theo từng phương tiện và tra cứu lịch sử nhận dạng kèm ảnh chứng cứ. Việc đóng gói pipeline AI trong tầng nghiệp vụ của Web App, kết hợp cơ chế xử lý bất đồng bộ và lưu trữ kết quả, cho phép kiểm chứng trực quan hoạt động của từng phiên nhận dạng thay vì chỉ dừng ở các mô hình thử nghiệm độc lập.

Bên cạnh những kết quả tích cực, đồ án cũng xác định rõ một số hạn chế còn tồn đọng như: hiện tượng trộn quỹ đạo (ID Switch) khi các phương tiện di chuyển quá gần nhau, và giới hạn của mô hình khi đối mặt với các biển số bị biến dạng vật lý hoặc che khuất hoặc mờ nghiêm trọng. Đây sẽ là tiền đề quan trọng cho các hướng phát triển ở giai đoạn tiếp theo.

\section{Hướng phát triển}

Từ những hạn chế còn tồn tại, hệ thống có thể được tiếp tục hoàn thiện và mở rộng theo các hướng nghiên cứu sau:
\begin{itemize}
    \item \textbf{Cải thiện thuật toán theo vết (Tracking):} Tích hợp các đặc trưng nhận dạng lại (ReID) chuyên biệt cho phương tiện giao thông để khắc phục hiện tượng trộn quỹ đạo trong môi trường mật độ giao thông cao.
    \item \textbf{Tăng cường xử lý đa khung hình:} Ứng dụng các thuật toán căn chỉnh và tổng hợp thông tin từ nhiều khung hình lân cận nhằm hỗ trợ khôi phục các biển số bị mờ nhòe do di chuyển tốc độ cao, kết hợp kiểm duyệt chặt chẽ để không làm sai lệch thông tin gốc.
    \item \textbf{Mở rộng tập dữ liệu và chuẩn hóa đánh giá:} Xây dựng bộ dữ liệu với đa dạng điều kiện thời tiết, góc máy, và loại camera; đồng thời thiết lập bộ tiêu chuẩn đánh giá (benchmark) độc lập cho từng module trong pipeline để đo lường chính xác hơn hiệu năng của từng công đoạn.
    \item \textbf{Tối ưu hóa nền tảng triển khai:} Chuyển đổi định dạng mô hình sang các nền tảng tối ưu cho thiết bị biên (như TensorRT, ONNX), áp dụng kỹ thuật lượng tử hóa (quantization) để giảm độ trễ, hướng tới việc tích hợp trực tiếp lên các camera giám sát thông minh (AI Camera).
\end{itemize}
\end{document}
